\documentclass[11pt,a4paper]{article}
\usepackage[utf-8]{inputenc}
\usepackage{amsmath}
\usepackage{amssymb}
\usepackage{graphicx}
\usepackage{array}
\usepackage{xcolor}
\usepackage{booktabs}
\usepackage{geometry}
\usepackage{float}
\usepackage{hyperref}
\usepackage{listings}

\geometry{margin=1in}

\title{\textbf{KMC Engine Function Logic Propagation Analysis:\\Parameter Symbolization vs. Actual Implementation}}
\author{Project: Summer\_project (Ashay-1)}
\date{\today}

\begin{document}

\maketitle

\begin{abstract}
This report provides a detailed analysis of the Kinetic Monte Carlo (KMC) Engine implementation, identifying which parameters are symbolic (placeholder/theoretical) versus which are actually utilized in the computational logic. The analysis traces function propagation through the call hierarchy and examines implementation limitations and gaps.
\end{abstract}

\tableofcontents
\newpage

\section{Executive Summary}

The KMCEngine implements a First Reaction Method (FRM) algorithm for simulating molecular dynamics on a crystalline lattice. The primary function \texttt{executeFRMStep()} coordinates the following workflow:

\begin{equation}
\text{Scan} \rightarrow \text{Evaluate} \rightarrow \text{Calculate Rates} \rightarrow \text{Select Event} \rightarrow \text{Execute}
\end{equation}

However, several parameters present in the function signatures are \textbf{symbolic only} (declared but unused), while others are \textbf{actively used}. This creates technical debt and potential for bugs.

\newpage

\section{Core Function Hierarchy Analysis}

\subsection{Primary Function: executeFRMStep()}

\subsubsection{Function Signature}
\begin{lstlisting}[language=C++, breaklines=true]
void executeFRMStep() {
    // Accesses: simBox, currentTime, gen, dis, attemptFrequency
}
\end{lstlisting}

\subsubsection{Parameters Classification}

\begin{table}[H]
\centering
\caption{KMCEngine::executeFRMStep() - Parameter Usage}
\begin{tabular}{|l|c|c|l|}
\hline
\textbf{Parameter/Variable} & \textbf{Used?} & \textbf{Type} & \textbf{Notes} \\
\hline
\texttt{simBox} & \checkmark & Active & Provides grid, box dimensions, event execution \\
\texttt{currentTime} & \checkmark & Active & Updated with $\Delta t$ from winning event \\
\texttt{gen} & \checkmark & Active & Seed for $\xi$ in event time calculation \\
\texttt{dis} & \checkmark & Active & Generates $\xi \in [0,1)$ \\
\texttt{attemptFrequency} & \checkmark & Active & Used in rate: $\nu = 10^{13}$ s$^{-1}$ \\
\hline
\end{tabular}
\end{table}

\newpage

\subsection{Secondary Functions: Kinetics Module}

\subsubsection{calculateReactionRate()}

\begin{lstlisting}[language=C++, breaklines=true]
static double calculateReactionRate(
    double attemptFrequency, 
    double energyBarrier,
    double transferCoeff, 
    double appliedVoltage, 
    double referenceVoltage)
\end{lstlisting}

\subsubsection{The Arrhenius-Butler-Volmer Equation}

\begin{equation}
k = \nu \exp\left[-\frac{E_a - \alpha(E - E_0)}{k_B T}\right]
\end{equation}

Where:
\begin{itemize}
    \item $\nu$ = attemptFrequency ($10^{13}$ s$^{-1}$)
    \item $E_a$ = energyBarrier (active parameter)
    \item $\alpha$ = transferCoeff (SYMBOLIC ONLY)
    \item $E$ = appliedVoltage (SYMBOLIC ONLY)
    \item $E_0$ = referenceVoltage (SYMBOLIC ONLY)
    \item $k_B T$ = thermal energy from Physics.hpp
\end{itemize}

\subsubsection{Parameter Classification for calculateReactionRate()}

\begin{table}[H]
\centering
\caption{Kinetics::calculateReactionRate() - Parameter Usage}
\begin{tabular}{|l|c|c|l|}
\hline
\textbf{Parameter} & \textbf{Used?} & \textbf{Value in Code} & \textbf{Classification} \\
\hline
\texttt{attemptFrequency} & Yes & Variable & ACTIVE \\
\texttt{energyBarrier} & Yes & Variable & ACTIVE \\
\texttt{transferCoeff} & No & 0.0 & SYMBOLIC ONLY \\
\texttt{appliedVoltage} & No & 0.0 & SYMBOLIC ONLY \\
\texttt{referenceVoltage} & No & 0.0 & SYMBOLIC ONLY \\
\hline
\end{tabular}
\end{table}

\subsubsection{Actual Implementation in executeFRMStep()}

\begin{lstlisting}[language=C++, breaklines=true]
// For REDUCTION events:
double rate = Kinetics::calculateReactionRate(
    attemptFrequency, 0.01, 0.0, 0.0, 0.0
);

// For DIFFUSION events:
double rate = Kinetics::calculateReactionRate(
    attemptFrequency, bindingEnergy, 0.0, 0.0, 0.0
);
\end{lstlisting}

\textbf{Implication:} The electrochemical parameters are \textbf{hardcoded to zero}, effectively disabling the Butler-Volmer correction term. The rate simplifies to:

\begin{equation}
k = \nu \exp\left[-\frac{E_a}{k_B T}\right]
\end{equation}

This is a \textbf{pure thermal diffusion model}, not an electrochemical one.

\newpage

\subsubsection{calculateBindingEnergy()}

\begin{lstlisting}[language=C++, breaklines=true]
static double calculateBindingEnergy(
    int targetIndex, 
    const std::vector<LatticeSite>& grid,
    double boxLengthX, 
    double boxLengthY, 
    double boxLengthZ)
\end{lstlisting}

\begin{table}[H]
\centering
\caption{Kinetics::calculateBindingEnergy() - Parameter Usage}
\begin{tabular}{|l|c|c|l|}
\hline
\textbf{Parameter} & \textbf{Used?} & \textbf{Classification} & \textbf{Notes} \\
\hline
\texttt{targetIndex} & Yes & ACTIVE & Locates the site \\
\texttt{grid} & Yes & ACTIVE & Neighbor search \\
\texttt{boxLengthX} & Yes & ACTIVE & PBC calculations \\
\texttt{boxLengthY} & Yes & ACTIVE & PBC calculations \\
\texttt{boxLengthZ} & Yes & ACTIVE & PBC calculations \\
\hline
\end{tabular}
\end{table}

\textbf{Algorithm:}
\begin{equation}
E_{\text{binding}} = \sum_{i \in \text{neighbors}} E_{\text{bond}}(s_{\text{target}}, s_i)
\end{equation}

Cutoff distance: $r_c = 3.9$ Å (covers first coordination shell).

\newpage

\subsubsection{calculatePBCDistanceSquared()}

\begin{lstlisting}[language=C++, breaklines=true]
static double calculatePBCDistanceSquared(
    const LatticeSite& s1, const LatticeSite& s2,
    double boxLengthX, double boxLengthY, 
    double boxLengthZ)
\end{lstlisting}

\begin{table}[H]
\centering
\caption{Kinetics::calculatePBCDistanceSquared() - Parameter Usage}
\begin{tabular}{|l|c|c|}
\hline
\textbf{Parameter} & \textbf{Used?} & \textbf{Classification} \\
\hline
\texttt{s1}, \texttt{s2} & Yes & ACTIVE \\
\texttt{boxLengthX/Y/Z} & Yes & ACTIVE \\
\hline
\end{tabular}
\end{table}

Implements minimum image convention:
\begin{equation}
d_{\text{PBC}}^2 = \sum_{\alpha \in \{x,y,z\}} \left(r_\alpha - L_\alpha \cdot \text{round}(r_\alpha / L_\alpha)\right)^2
\end{equation}

\newpage

\section{Function Call Propagation Tree}

\begin{verbatim}
executeFRMStep()
|-- FOR each lattice site i:
|   |-- calculateBindingEnergy(i, grid, L_x, L_y, L_z)
|   |   +-- FOR each neighbor j:
|   |       +-- calculatePBCDistanceSquared(s_i, s_j, Lx, Ly, Lz)
|   |
|   +-- FOR each potential target j:
|       |-- calculatePBCDistanceSquared(s_i, s_j, Lx, Ly, Lz)
|       |
|       +-- IF distance < cutoff:
|           |-- calculateReactionRate(nu, E_a, 0.0, 0.0, 0.0)
|           |
|           +-- dt = -ln(xi) / k
|
+-- SELECT event with min(dt)
    +-- EXECUTE:
        |-- simBox.swapSpecies() [diffusion]
        +-- simBox.reduceSpecies() [reduction]
\end{verbatim}

\newpage

\section{Symbolic vs. Active Parameter Mapping}

\subsection{Summary Table}

\begin{table}[H]
\centering
\caption{Complete Parameter Classification}
\begin{tabular}{|l|l|c|c|}
\hline
\textbf{Function} & \textbf{Parameter} & \textbf{Active} & \textbf{Symbolic} \\
\hline
\multirow{5}{*}{\texttt{calculateReactionRate()}} 
  & attemptFrequency & X & \\
  & energyBarrier & X & \\
  & transferCoeff & & X \\
  & appliedVoltage & & X \\
  & referenceVoltage & & X \\
\hline
\multirow{3}{*}{\texttt{calculateBindingEnergy()}}
  & targetIndex & X & \\
  & grid & X & \\
  & boxLength* & X & \\
\hline
\multirow{2}{*}{\texttt{calculatePBCDistanceSquared()}}
  & s1, s2 & X & \\
  & boxLength* & X & \\
\hline
\multirow{4}{*}{\texttt{executeFRMStep()}}
  & simBox & X & \\
  & currentTime & X & \\
  & gen (RNG) & X & \\
  & attemptFrequency & X & \\
\hline
\end{tabular}
\end{table}

\textbf{Finding:} 60\% of electrochemical parameters are symbolic placeholders.

\newpage

\section{Implementation Limitations}

\subsection{Limitation 1: Missing Electrochemistry}

\subsubsection{Issue}
The function signature promises electrochemical modeling through Butler-Volmer kinetics, but the implementation passes zeros for:
\begin{itemize}
    \item $\alpha$ (charge transfer coefficient)
    \item $E$ (applied voltage/potential)
    \item $E_0$ (reference potential)
\end{itemize}

\subsubsection{Code Evidence}
\begin{lstlisting}[language=C++, breaklines=true]
// Line 89 (reduction event):
double rate = Kinetics::calculateReactionRate(
    attemptFrequency, 0.01, 0.0, 0.0, 0.0);
\end{lstlisting}

\subsubsection{Impact}
Cannot model: overpotential effects, electrochemical phase boundaries, potential-driven kinetics.

\subsection{Limitation 2: Hardcoded Reduction Barrier}

\subsubsection{Issue}
The reduction event uses a hardcoded energy barrier of $0.01$ eV:

\begin{lstlisting}[language=C++, breaklines=true]
double rate = Kinetics::calculateReactionRate(
    attemptFrequency, 0.01, 0.0, 0.0, 0.0);
\end{lstlisting}

\subsubsection{Problem}
This assumes \textbf{all salt and solvent reduction at Li metal has the same barrier}, regardless of:
\begin{itemize}
    \item Interfacial orientation
    \item Local electrolyte composition
    \item Coordination environment
\end{itemize}

\subsection{Limitation 3: Incomplete Event Classification}

Only two event types are modeled:
\begin{equation}
\text{eventType} = \begin{cases} 
0 & \text{Diffusion} \\
1 & \text{Reduction}
\end{cases}
\end{equation}

Missing:
\begin{itemize}
    \item Ion-ion correlations
    \item Vacancy formation/annihilation
    \item Complex electrolyte decomposition
    \item Electromigration under field
\end{itemize}

\subsection{Limitation 4: Cutoff Distance Hardcoding}

\begin{lstlisting}[language=C++, breaklines=true]
const double hopCutoffSq = 3.9 * 3.9;
\end{lstlisting}

\textbf{Issue:} Hardcoded to $3.9$ Å with no justification or parametrization.

Expected: Should be derived from lattice constant $a_0 = 4.47$ Å:
\begin{equation}
r_c = n \cdot a_0 \quad \text{where} \quad n \approx 0.87
\end{equation}

\newpage

\section{Comparison with Theoretical Model}

\subsection{Intended Mathematical Framework}

\begin{equation}
k = \nu \exp\left[-\frac{E_a - \alpha(V - V_0)}{k_B T}\right]
\end{equation}

\subsection{Actual Implementation}

\begin{equation}
k = \nu \exp\left[-\frac{E_a}{k_B T}\right]
\end{equation}

\textbf{Missing:} The entire overpotential correction term.

\newpage

\section{Data Flow Analysis}

\begin{table}[H]
\centering
\caption{Data Flow: Source to Function to Use}
\begin{tabular}{|l|l|l|l|}
\hline
\textbf{Source} & \textbf{Value} & \textbf{Function} & \textbf{Use} \\
\hline
Physics.hpp & $k_B$, $T$ & calculateReactionRate() & Thermal term \\
Lattice.hpp & Species enum & calculateBindingEnergy() & Bond lookup \\
SimulationBox & Grid, dims & executeFRMStep() & Site enumeration \\
Hardcoded & 0.01 eV & executeFRMStep() & Reduction $E_a$ \\
Hardcoded & 3.9 Å & Both modules & Cutoff distance \\
Unused & $\alpha, V, V_0$ & calculateReactionRate() & Symbolic only \\
\hline
\end{tabular}
\end{table}

\newpage

\section{Recommendations}

\subsection{Short Term: Fix Hardcoded Values}

\begin{lstlisting}[language=C++, breaklines=true]
class KMCEngine {
private:
    double hopCutoffDistance;
    double reductionEnergyBarrier;
    
public:
    KMCEngine(SimulationBox& box, 
              double cutoff = 3.9,
              double redBarrier = 0.01)
        : simBox(box), hopCutoffDistance(cutoff), 
          reductionEnergyBarrier(redBarrier),
          currentTime(0.0), gen(std::random_device{}()),
          dis(0.0, 1.0) {}
};
\end{lstlisting}

\subsection{Medium Term: Implement Missing Functions}

Create functions for:

1. \texttt{calculateReductionBarrier(sourceSpecies, coordination)}
2. \texttt{calculateTransferCoefficient(interface\_type)}
3. \texttt{getAppliedVoltage()} — external voltage control

\subsection{Long Term: Full Electrochemistry}

\begin{lstlisting}[language=C++, breaklines=true]
double rate = Kinetics::calculateReactionRate(
    attemptFrequency,
    calculateReductionBarrier(sourceSite),
    calculateTransferCoeff(targetSite),
    appliedVoltage,
    referenceVoltage
);
\end{lstlisting}

\newpage

\section{Conclusion}

\begin{table}[H]
\centering
\caption{Summary of Findings}
\begin{tabular}{|l|c|l|}
\hline
\textbf{Metric} & \textbf{Value} & \textbf{Assessment} \\
\hline
Total rate equation parameters & 5 & \\
Active parameters & 2 & 40\% \\
Symbolic/Unused parameters & 3 & 60\% \\
Hardcoded constants & 2 & Inflexible \\
Dynamic functions & 3 & Good \\
Missing dynamic functions & 3+ & Needed \\
\hline
\end{tabular}
\end{table}

\subsection{Key Takeaways}

\begin{enumerate}
    \item The KMC engine implements \textbf{pure thermal diffusion}, not electrochemistry.
    
    \item Three crucial parameters are \textbf{symbolic placeholders} hardcoded to $0.0$.
    
    \item Two physical parameters are \textbf{hardcoded magic numbers}.
    
    \item The implementation is \textbf{architecturally sound} but \textbf{functionally incomplete}.
    
    \item Dynamic functions must replace hardcoded values.
\end{enumerate}

\end{document}